\section{Noise sensitivity and agnostic learning}\label{sec:consequences}

The influence bound has two consequences that are independent of the
weighted grading used to prove it. A random-bucketing construction
transfers the bound from one-bit changes to independent noise.
Smoothing by this noise then gives low-degree polynomial approximants,
which yield an agnostic learner by $L_1$ regression.

\subsection{Noise sensitivity}

For $0\le\eta\le1/2$, define the noise sensitivity
$\operatorname{NS}_{\eta}(f)=\Prb\{f(X)\ne f(Y)\}$, where $X$ is uniform
on the cube and $Y$ is obtained by flipping each coordinate of $X$
independently with probability $\eta$.

\begin{corollary}[Boolean noise sensitivity]\label{cor:noise-sensitivity}
Let $n\ge1$ and $d\ge0$ be integers. If $f=\sgn(p)$ on
$\{-1,1\}^n$ for a real polynomial $p$ of degree at most $d$, with
$\sgn(0)=1$, then
\[
 \operatorname{NS}_{\eta}(f)\le C d\sqrt\eta
 \qquad(0<\eta\le1/2),
\]
where $C$ is an absolute constant.
\end{corollary}

\begin{proof}
Constants have zero noise sensitivity. For an integer $q\ge2$, use the
random-bucketing reduction of Diakonikolas, Raghavendra, Servedio, and
Tan~\cite[Section~7.1, equation~(14)]{DRST2014}: choose independent uniform
signs $a_i$ and independent uniform buckets $b(i)\in\{1,\ldots,q\}$,
and set
\[
 g(z)=f(a_1z_{b(1)},\ldots,a_nz_{b(n)}),
 \qquad z\in\{-1,1\}^q.
\]
Multilinearizing the substituted polynomial preserves its cube values,
including zeros, and gives degree at most $\min(d,q)$.
Theorem~\ref{thm:main} therefore gives $I(g)\le8d\sqrt q$; a constant
substitution contributes zero.

Independently choose uniform $z\in\{-1,1\}^q$ and $J\in\{1,\ldots,q\}$.
Put $X_i=a_i z_{b(i)}$ and
$Y_i=X_i(-1)^{\mathbf1_{\{b(i)=J\}}}$.
Conditional on $J$, the flip indicators are independent Bernoulli
variables of parameter $1/q$, and their joint law does not depend on $J$.
Conditional on $b,J,z$, the random signs make $X$ uniform, so $X$ is
independent of the flip vector. Thus $(X,Y)$ has exactly the
noise law at rate $1/q$. Since replacing $z$ by $z^{\oplus J}$ produces
$Y$, averaging first over $z,J$ gives
\[
 \operatorname{NS}_{1/q}(f)=\frac{\E_{a,b}I(g)}q
 \le\frac{8d}{\sqrt q}.
\]

For the rest of this section, Fourier coefficients and $L_1,L_2$ norms
are normalized by uniform probability on the cube. In particular,
$\widehat f(S)=\E[f(X)\chi_S(X)]$ and
\[
 \operatorname{NS}_{\eta}(f)
 =\frac12\sum_{S\subseteq[n]}
 \bigl(1-(1-2\eta)^{|S|}\bigr)\widehat f(S)^2.
\]
This expression is nondecreasing for $0\le\eta\le1/2$.
Taking $q=\lfloor1/\eta\rfloor$, so that
$\eta\le1/q\le1/2$ and $q\ge1/(2\eta)$, proves the claim with
$C=8\sqrt2$.
\end{proof}

For comparison, Kane proved an absolute $O(d\sqrt\varepsilon)$ noise
sensitivity bound for degree-$d$ polynomial threshold functions under
standard Gaussian noise~\cite[Theorem~1]{Kane2011}.
Here $0<\varepsilon\le1$, and the input pair is
$X$ and $(1-\varepsilon)X+\sqrt{2\varepsilon-\varepsilon^2}\,Z$,
with $X,Z$ independent standard Gaussian vectors.
His proof counts sign changes along Gaussian rotations.
Corollary~\ref{cor:noise-sensitivity} gives the same dependence on degree
and noise rate for independent coordinate flips on the Boolean cube.

\subsection{Agnostic learning}

In agnostic learning the labels may be arbitrary: the objective is to
predict nearly as well as the best function in a specified class.
We use the $L_1$ polynomial regression method of Kalai, Klivans,
Mansour, and Servedio~\cite{KKMS2008}, which was applied to polynomial
threshold functions in~\cite[Section~8]{DRST2014}.

\begin{corollary}[Uniform-marginal agnostic learning]\label{cor:agnostic-learning}
Let $n,d\ge1$ be integers, and let $\mathcal C_{n,d}$ be the class of
functions $\sgn(p)$ on $\{-1,1\}^n$ with $\deg p\le d$.
Let $D$ be any distribution on $\{-1,1\}^n\times\{-1,1\}$ whose first
marginal is uniform, and set
\[
 \mathrm{OPT}=\min_{f\in\mathcal C_{n,d}}\Prb_{(X,Y)\sim D}\{f(X)\ne Y\}.
\]
There is an absolute constant $A$ such that, for $0<\alpha,\delta<1$,
a learner using independent labeled samples from $D$ returns, with
probability at least $1-\delta$, a Boolean hypothesis $h$ satisfying
\[
 \Prb_{(X,Y)\sim D}\{h(X)\ne Y\}\le\mathrm{OPT}+\alpha.
\]
Its running time and sample size are polynomial in
$(n+1)^k$, $1/\alpha$, and $\log(1/\delta)$, where
\[
 k=\min\left\{n,\left\lceil A d^2\alpha^{-2}\log(4/\alpha)\right\rceil\right\}.
\]
The labels are unrestricted, and $h$ need not belong to $\mathcal C_{n,d}$.
\end{corollary}

\begin{proof}
We first construct an $L_1$ approximant to each $f\in\mathcal C_{n,d}$.
For $0<\eta<1/2$ and $\rho=1-2\eta$, the noise operator is
\[
 T_\rho f=\sum_{S\subseteq[n]}\rho^{|S|}\widehat f(S)\chi_S.
\]
Equivalently, $T_\rho f(x)$ is the expected value of $f$ after independent
noise of rate $\eta$ is applied to $x$. Since $f$ is Boolean and
$T_\rho f$ takes values in $[-1,1]$,
$|f-T_\rho f|=1-fT_\rho f$ pointwise. Corollary~\ref{cor:noise-sensitivity}
therefore gives, with $C=8\sqrt2$,
\[
 \norm{f-T_\rho f}_1=2\operatorname{NS}_\eta(f)\le2Cd\sqrt\eta.
\]
Take $\eta=(\alpha/(8Cd))^2$ and
$K=\lceil\log(4/\alpha)/(2\eta)\rceil$.
Let $q$ be the Fourier truncation of $T_\rho f$ to degrees at most
$\min(K,n)$. If $K<n$, Parseval's identity gives
\[
 \norm{T_\rho f-q}_1\le\norm{T_\rho f-q}_2
 \le\rho^{K+1}\le e^{-2\eta(K+1)}\le\alpha/4;
\]
if $K\ge n$, the truncation error is zero.
The first error is also at most $\alpha/4$, so
\begin{equation}\label{eq:learning-approximation}
 \norm{f-q}_1\le\alpha/2.
\end{equation}
Taking $A=32C^2$ makes $\min(K,n)=k$ in the statement.

For completeness, we explain how this direct $L_1$ estimate enters the
regression argument of~\cite[proof of Theorem~5]{KKMS2008}.
Use the $M=\sum_{j=0}^k\binom nj\le(n+1)^k$ characters of degree at
most $k$ as features. On a training sample $Z$ of size $m$, fit a
polynomial $P$ in these features by minimizing empirical absolute loss,
then choose $t\in[-1,1]$ to minimize the empirical error of
$h=\sgn(P-t)$. The fit is a linear program, and the threshold is found by
sorting the fitted values. Write $\widehat{\E}_Z$ and
$\widehat{\operatorname{err}}_Z$ for empirical means and errors.
A uniformly random threshold has error at most half the absolute loss,
so the minimizing threshold satisfies
\[
 \widehat{\operatorname{err}}_Z(h)
 \le\tfrac12\widehat{\E}_Z|Y-P(X)|
 \le\tfrac12\widehat{\E}_Z|Y-q(X)|.
\]
Choose an optimal $f\in\mathcal C_{n,d}$ and its fixed approximant $q$.
The triangle inequality and \eqref{eq:learning-approximation} imply
$\E_Z\widehat{\operatorname{err}}_Z(h)\le\mathrm{OPT}+\alpha/4$.
These hypotheses are halfspaces in the $M$ character features, so their
VC dimension is at most $M+1$.
The usual expected uniform generalization bound, as used
in~\cite[proof of Theorem~5]{KKMS2008}, makes
$\E_Z\operatorname{err}_D(h)\le\mathrm{OPT}+3\alpha/8$
for $m$ polynomial in $M$ and $1/\alpha$.

Apply Markov's inequality to the nonnegative random variable
$\operatorname{err}_D(h)$. Since $\mathrm{OPT}\le1$ and $\alpha<1$,
\[
 \Prb_Z\{\operatorname{err}_D(h)>\mathrm{OPT}+\alpha/2\}
 \le\frac{\mathrm{OPT}+3\alpha/8}{\mathrm{OPT}+\alpha/2}
 \le1-\alpha/12.
\]
Thus one trial succeeds with probability at least $\alpha/12$.
Repeat independently
$r=\lceil(12/\alpha)\log(2/\delta)\rceil$ times.
With probability at least $1-\delta/2$, one candidate has this error.
On an independent validation sample of size
$\lceil8\alpha^{-2}\log(4r/\delta)\rceil$, Hoeffding's inequality and a
union bound estimate all $r$ errors within $\alpha/4$, except with
probability $\delta/2$. Selecting the best validation score therefore
gives error at most $\mathrm{OPT}+\alpha$. All the training,
regression, repetition and validation costs have the stated polynomial
bounds.
\end{proof}
